% arXiv build only (main-arxiv.tex). Sources: docs/theory/proposition1.md;
% results-hust.md, results-uored.md, results-ptbxl.md (earlier code).
\section{Proof of Proposition~\ref{prop:anyselector}}
\label{app:proof}

\paragraph{Step 1: the loss is a threshold function.} Let $m(x) = \min\{\gh_r(x) : d_r(x) = 1\}$,
with $m(x) = +\infty$ if no representation disagrees. Because $r_0$ never disagrees, a member of
$\Sset(x)$ disagrees iff some $r$ has $d_r(x) = 1$ and $\gh_r(x) \le \lambda$, so
$L(x;\lambda) = \1[\lambda \ge m(x)]$. Hence $\lambda \mapsto L(x;\lambda)$ is non-decreasing,
bounded by 1, and $L(x;-1) = 0$ because $\gh \ge 0 > -1$.

\paragraph{Step 2: exchangeable losses.} $m$ is a fixed measurable function of $x$ once the
predictor and the heads are fixed, and both were fitted on data independent of
$(x_1,\dots,x_{n+1})$. Hence $L_i(\lambda) := L(x_i;\lambda)$, $i = 1,\dots,n+1$, are exchangeable
random functions.

\paragraph{Step 3: conformal risk control on a finite grid.} This is the argument of Theorem~1 of
\cite{angelopoulos2024crc}, written for a finite grid $\Lambda$ so that no continuity of the loss
is needed: with $\Sset = \{\gh \le \lambda\}$ the loss $\1[\lambda \ge m]$ is not right-continuous
in the direction their theorem assumes, and a supremum over a continuum would not be attained when
scores tie. Let $\hat R_{n+1}(\lambda) = \frac{1}{n+1}\sum_{i\le n+1} L_i(\lambda)$ and
$\lh' = \max\{\lambda \in \Lambda : \hat R_{n+1}(\lambda) \le \alpha\}$. The set contains $-1$,
where every loss is 0, and $\Lambda$ is finite, so the maximum is attained and
$\hat R_{n+1}(\lh') \le \alpha$.
\begin{itemize}
  \item Because $L_{n+1} \le 1$, $\hat R_{n+1}(\lambda) \le (n\hat R_n(\lambda) + 1)/(n+1)$ for
  every $\lambda$. Every $\lambda$ that qualifies for $\lh$ in \eqref{eq:lhat} therefore qualifies
  for $\lh'$, so $\lh \le \lh'$ (also when $\lh = -1$), and by monotonicity
  $L_{n+1}(\lh) \le L_{n+1}(\lh')$.
  \item $\lh'$ is a symmetric function of the exchangeable $L_1,\dots,L_{n+1}$, hence
  $\E[L_{n+1}(\lh')] = \E\bigl[\frac{1}{n+1}\sum_i L_i(\lh')\bigr] = \E[\hat R_{n+1}(\lh')] \le
  \alpha$.
\end{itemize}
Together, $\E[L(x_{n+1};\lh)] \le \alpha$.

\paragraph{Step 4: any selector inside the safe set.} On $C$, $\pi(x_{n+1},z) \in
S_{\lh}(x_{n+1})$, therefore $d_\pi(x_{n+1})\1\{C\} \le \max_{r\in S_{\lh}} d_r(x_{n+1}) =
L(x_{n+1};\lh)$ pointwise, whatever the selector looks at: scores, system state, its own
randomness, even the losses themselves. Taking expectations gives the first inequality. What must
not depend on the test losses is $\lh$, the scores $\gh$ (both fixed before the test window), and
the set of windows the average is taken over. \qed

\paragraph{Computing $\lh$.} With $m_{(1)} \le \dots \le m_{(n)}$ the sorted calibration values and
$k = \lfloor \alpha(n+1) - 1\rfloor$: if $k < 0$ (that is, $n < \lceil 1/\alpha\rceil - 1$), even
zero observed losses fail the bound and the safe set is raw only; if $k \ge n$ or $m_{(k+1)} =
+\infty$, everything is safe; otherwise $\lh$ is the largest float strictly below $m_{(k+1)}$, so
ties at $m_{(k+1)}$ are excluded.

\paragraph{The cascade.} If $F$ and $A(x)$ are each calibrated by V0 at $\alpha/2$ (each a fixed
score function fitted on \texttt{ptrain}, each threshold from \texttt{calib}), then for every window
$\max_{F\cup A(x)} d \le \max_F d + \max_{A(x)} d$, hence $\E[\max_{F\cup A(x)} d] \le \alpha$, and
Step~4 covers any selector inside the union, including one that decides from the system state
whether to compute $\gh$, as long as it sends a member of $F$ when it does not.

\section{Datasets run with earlier code}
\label{app:earlier}

% src: results-hust.md, results-uored.md, results-ptbxl.md (E1, E2, E3, E5, E7)
HUST, UORED and PTB-XL were run before the code of the final rerun (non-blocking pacing,
lag-aware deadlines, audit feedback for certified decisions only, group CRC). Their numbers are
not comparable with Section~\ref{sec:results} number for number; the HUST run predates the
cascade. Table~\ref{tab:earlier} reports them at $\alpha = 0.05$.

\begin{table}[h]
\centering
\caption{Earlier-code datasets, five folds each, $\alpha = 0.05$.}
\label{tab:earlier}
\small
\begin{tabular}{lrrr}
\toprule
& HUST & UORED & PTB-XL \\
\midrule
$\fraw$ macro-F1 (majority)     & 0.636 (0.167) & 0.751 (0.267) & 0.673 (0.261) \\
V0 risk                          & 0.0777 & 0.0788 & 0.0489 \\
Cascade risk\expl                & -- & 0.0588 & 0.0244 \\
Static risk\expl                 & 0.0265 & 0.0173 & 0.0320 \\
H3 cost vs B4                    & $-27.2\,\%$ & $-15.7\,\%$ & $-61.4\,\%$ \\
H5 full menu vs best family      & $+15.1\,\%$ (M) & $+15.2\,\%$ (M) & $+15.9\,\%$ (S) \\
H7 worst p50 gap                 & 11.9\,\% & 13.1\,\% & 23.4\,\% \\
\bottomrule
\end{tabular}
\end{table}

On HUST and UORED, V0's mean risk is above $\alpha$ at every level, and on UORED the cascade's mean
is above $\alpha$ at $\alpha = 0.05$. H7 fails on all three. These runs predate the non-blocking
pacing fix, which changed the Speech Commands H7 result from a 17.3\,\% gap and 84.0\,\% agreement to
12.5\,\% and 98.7\,\%; they were not rerun.
